% Lesson 2 theme: Lesson 1 preamble, then three overrides.
% =============================================================
% Agentic AI for Economic Research, Lesson 1 theme (v2)
% Compiler: pdflatex (no fontspec). Paths are relative to repo root.
% Design read: lecture deck for economics PhD students and faculty,
% watched over a screen share; calm academic language; one ink colour,
% one accent, real screenshots and data charts as the only visuals.
% Dials: variance 5, motion 1, density 4.
% =============================================================

\documentclass[aspectratio=169,10pt]{beamer}

\usepackage[T1]{fontenc}
\usepackage[utf8]{inputenc}
\usepackage[english]{babel}

\usepackage[default,scale=0.95]{sourcesanspro}
\usepackage[scaled=0.85]{sourcecodepro}
\renewcommand{\familydefault}{\sfdefault}
\usepackage{microtype}

\usepackage{graphicx}
\usepackage{xcolor}
\newcommand{\rev}[1]{#1}
\usepackage{tikz}
\usetikzlibrary{positioning,arrows.meta,calc,fit,backgrounds}
\usepackage{pgfplots}
\usepackage{pgfplotstable}
\pgfplotsset{compat=1.18}
\usepgfplotslibrary{dateplot}
\usepackage{tabularx}
\usepackage{booktabs}
\usepackage{array}
\usepackage{ragged2e}
\usepackage{hyperref}
\usepackage[most]{tcolorbox}
\usepackage{fancyvrb}
\usepackage[round,authoryear]{natbib}
\bibliographystyle{plainnat}

\ifdefined\DeclareUnicodeCharacter
  \DeclareUnicodeCharacter{2192}{\ensuremath{\rightarrow}}
  \DeclareUnicodeCharacter{2190}{\ensuremath{\leftarrow}}
  \DeclareUnicodeCharacter{2717}{x}
  \DeclareUnicodeCharacter{00B7}{\textperiodcentered}
\fi

\graphicspath{{assets/lesson1/screens/}{assets/lesson1/logos/}{assets/covers/}{assets/real/}}

% One cover. Same band, same fonts, same lines. \coverframe{image}{title}{subtitle}{n}
\newcommand{\coverframe}[4]{%
\begin{frame}[plain]
  \begin{tikzpicture}[remember picture,overlay]
    \node[anchor=center,inner sep=0] at (current page.center)
      {\includegraphics[width=\paperwidth,height=\paperheight]{#1}};
    \fill[black,opacity=.55] (current page.south west) rectangle ($(current page.south west)+(\paperwidth,3.45cm)$);
    \node[anchor=south west,text width=13.4cm,align=left,inner sep=0pt] at ($(current page.south west)+(0.7cm,0.38cm)$)
      {{\color{white}\fontsize{25}{29}\selectfont\bfseries #2}\\[4pt]
       {\normalsize\color{white!85} #3}\\[5pt]
       {\footnotesize\color{white!78} Agentic AI for Economic Research \quad|\quad Lesson #4 \quad|\quad Marco Le Moglie}\\[4pt]
       {\hypersetup{urlcolor=white!85}\href{https://github.com/MarcoLeMoglie/Agentic_Economic_Research}{\footnotesize\color{white!85}\includegraphics[height=1.05em]{github_mark_white.png}\ GitHub repo}}};
  \end{tikzpicture}
\end{frame}}

% ---------- Palette (locked: one ink, one accent, one cool grey family) ----------
\definecolor{ink}{HTML}{163A59}
\definecolor{accent}{HTML}{D97706}
\definecolor{dark}{HTML}{111827}
\definecolor{muted}{HTML}{4B5563}
\definecolor{faint}{HTML}{9CA3AF}
\definecolor{hair}{HTML}{D1D5DB}
\definecolor{paper}{HTML}{F8FAFC}
\definecolor{tint}{HTML}{EEF2F6}
% semantic only (good / bad), never decorative
\definecolor{good}{HTML}{166534}
\definecolor{bad}{HTML}{B91C1C}
\definecolor{stepR}{HTML}{6D28D9}
\definecolor{stepE}{HTML}{0F766E}

\hypersetup{colorlinks=true,linkcolor=ink,urlcolor=ink,citecolor=ink}

% ---------- Beamer base ----------
\useinnertheme{rectangles}
\useoutertheme{default}
\setbeamersize{text margin left=2.2em,text margin right=2.2em}
\setbeamercolor{background canvas}{bg=paper}
\setbeamercolor{normal text}{fg=dark}
\setbeamercolor{structure}{fg=ink}
\setbeamercolor{alerted text}{fg=accent}
\setbeamertemplate{navigation symbols}{}

\setbeamerfont{frametitle}{size=\Large,series=\bfseries}
\setbeamercolor{frametitle}{fg=ink}
\setbeamertemplate{frametitle}{%
  \vspace{0.9em}%
  \begin{minipage}[t]{\dimexpr\paperwidth-4.4em\relax}%
    \raggedright\insertframetitle\par
    \vspace{-0.35em}{\color{accent}\rule{2.2em}{1.2pt}}%
  \end{minipage}%
  \vspace{0.2em}%
}

\setbeamertemplate{footline}{%
  \leavevmode\hbox to \paperwidth{%
    \hspace{2.2em}{\color{faint}\fontsize{6.5}{8}\selectfont Agentic AI for Economic Research, Lesson 1}%
    \hfill{\color{faint}\fontsize{6.5}{8}\selectfont\insertframenumber}\hspace{2.2em}}%
  \vspace{0.9em}%
}

% bullets: small grey square, no chevrons
\setbeamertemplate{itemize item}{\raisebox{0.18em}{\color{faint}\rule{0.32em}{0.32em}}}
\setbeamertemplate{itemize subitem}{\raisebox{0.2em}{\color{hair}\rule{0.26em}{0.26em}}}
\setbeamertemplate{enumerate item}{\color{ink}\bfseries\insertenumlabel.}
\setlength{\leftmargini}{1.1em}

% ---------- tcolorbox styles (radius locked at 2pt) ----------
\tcbset{
  l1base/.style={enhanced,arc=2pt,outer arc=2pt,boxrule=0pt,
    left=0pt,right=0pt,top=0pt,bottom=0pt},
}

% Definition: accent rule on the left, term, text, source
\newtcolorbox{defbox}[1]{enhanced,frame hidden,interior hidden,
  borderline west={1.6pt}{0pt}{accent},
  left=0.9em,right=0pt,top=0.15em,bottom=0.15em,
  fontupper=\normalsize,
  before upper={{\color{ink}\bfseries #1}\par\vspace{0.25em}}}
\newcommand{\defn}[3]{%
  \begin{defbox}{#1}#2\if\relax\detokenize{#3}\relax\else\par\vspace{0.3em}{\color{muted}\scriptsize #3}\fi\end{defbox}}

% Real file or real terminal output, read verbatim from disk:
% \filebox[width]{path}{caption}
\newcommand{\filebox}[3][\linewidth]{%
  \begin{tcolorbox}[l1base,width=#1,colback=white,
      borderline={0.4pt}{0pt}{hair},left=0.6em,right=0.4em,top=0.4em,bottom=0.4em]
    \VerbatimInput[fontsize=\fontsize{6}{7.1}\selectfont,formatcom=\color{dark}]{#2}%
  \end{tcolorbox}%
  \par\vspace{0.15em}{\color{muted}\scriptsize #3}}

% Screenshot frame: hairline border, caption below
\newcommand{\shot}[3][\linewidth]{%
  \begin{tcolorbox}[l1base,width=#1,colback=white,
      borderline={0.4pt}{0pt}{hair},left=0pt,right=0pt,top=0pt,bottom=0pt]
    \includegraphics[width=\linewidth]{#2}%
  \end{tcolorbox}%
  \par\vspace{0.15em}{\color{muted}\scriptsize #3}}
\newcommand{\shoth}[4][\linewidth]{% height-capped variant: [width]{file}{height}{caption}
  \begin{tcolorbox}[l1base,width=#1,colback=white,halign=center,
      borderline={0.4pt}{0pt}{hair},left=0pt,right=0pt,top=0pt,bottom=0pt]
    \includegraphics[width=\linewidth,height=#3,keepaspectratio]{#2}%
  \end{tcolorbox}%
  \par\vspace{0.15em}{\color{muted}\scriptsize #4}}

% Source line at the bottom of a frame
\newcommand{\source}[1]{\vfill{\color{muted}\scriptsize Source: #1\par}}

% One statement per slide, at most
\newcommand{\takeaway}[1]{\par\vspace{0.5em}{\color{ink}\large\bfseries #1\par}}

% Small group label used instead of numbered labels
\newcommand{\group}[1]{{\color{ink}\bfseries #1}\par\vspace{0.2em}}

% Section divider: ink background, topic name only
\newcommand{\sectionframe}[2][]{%
\begin{frame}[plain]
  \begin{tikzpicture}[remember picture,overlay]
    \fill[ink] (current page.south west) rectangle (current page.north east);
    \node[anchor=south west,inner sep=0pt,text=white,align=left,text width=0.8\paperwidth]
      at ($(current page.south west)+(1.6cm,1.7cm)$)
      {{\fontsize{28}{33}\selectfont\bfseries #2}\par\vspace{0.4em}{\color{white!75!ink}\large #1}};
    \fill[accent] ($(current page.south west)+(1.6cm,1.25cm)$) rectangle ++(1.4cm,0.06cm);
  \end{tikzpicture}
\end{frame}}

% Live demo placeholder
\newcommand{\demoframe}[3]{% {title}{what we do}{what to look at}
\begin{frame}{#1}
  \begin{columns}[T,onlytextwidth]
    \column{0.55\textwidth}
      \group{What we do}
      #2
    \column{0.40\textwidth}
      \group{What to look at}
      #3
  \end{columns}
\end{frame}}

% ---------- TikZ styles ----------
\tikzset{
  lab/.style={align=center,font=\small,text=dark},
  sub/.style={align=center,font=\scriptsize,text=muted},
  hl/.style={align=center,font=\small\bfseries,text=ink},
  frameb/.style={draw=hair,line width=0.5pt,rounded corners=2pt,fill=white},
  inkb/.style={draw=ink,line width=0.8pt,rounded corners=2pt,fill=white},
  accb/.style={draw=accent,line width=0.9pt,rounded corners=2pt,fill=white},
  arr/.style={-{Stealth[length=4pt]},line width=0.6pt,color=faint},
  arrA/.style={-{Stealth[length=5pt]},line width=0.9pt,color=accent},
}

% ---------- pgfplots defaults: no legends, no background tracks ----------
\pgfplotsset{
  l1axis/.style={
    axis lines*=left,
    axis line style={hair},
    tick style={draw=none},
    ticklabel style={font=\scriptsize,color=muted},
    label style={font=\scriptsize,color=muted},
    every axis plot/.append style={line width=0.6pt},
    enlargelimits=false,
    clip=false,
    /pgf/number format/assume math mode=true,
  },
  % keep only rows whose column #1 equals #2 (colour bars by group from the CSV)
  only if/.style 2 args={
    x filter/.code={\getthisrow{#1}\tempa\edef\tempb{#2}%
      \ifx\tempa\tempb\else\def\pgfmathresult{inf}\fi}},
}

% ---------- Metadata ----------
\title{Agentic AI for Economic Research}
\author{Marco Le Moglie}
\institute{Catholic University of the Sacred Heart}


\setbeamertemplate{footline}{%
  \leavevmode\hbox to \paperwidth{%
    \hfill{\color{faint}\fontsize{6.5}{8}\selectfont\insertframenumber}\hspace{2.2em}}%
  \vspace{0.9em}%
}

\graphicspath{{assets/lesson2/screens/}{assets/lesson2/illustrations/}{assets/lesson2/diagrams/}{assets/lesson1/logos/}{assets/real/}}

% Macros used by the Lesson 2 source pack. Colours and frame chrome stay on the Lesson 1 theme.
\ifdefined\DeclareUnicodeCharacter
  \DeclareUnicodeCharacter{2013}{--}
  \DeclareUnicodeCharacter{2014}{---}
  \DeclareUnicodeCharacter{2018}{`}
  \DeclareUnicodeCharacter{2019}{'}
  \DeclareUnicodeCharacter{201C}{``}
  \DeclareUnicodeCharacter{201D}{''}
\fi
% Review colour accepted: \chg keeps the words and no longer paints them.
\newcommand{\chg}[1]{#1}
\newcommand{\chgshoth}[4][\linewidth]{%
  \begin{tcolorbox}[l1base,width=#1,colback=white,halign=center,
      borderline={1pt}{0pt}{hair},left=0pt,right=0pt,top=0pt,bottom=0pt]
    \includegraphics[width=\linewidth,height=#3,keepaspectratio]{#2}%
  \end{tcolorbox}%
  \par\vspace{0.15em}{\color{muted}\scriptsize #4}}
\newtcolorbox{chgfig}{l1base,colback=paper,borderline={1pt}{0pt}{hair},left=4pt,right=4pt,top=4pt,bottom=4pt,halign=center}

% YouTube clip card: \ytclip{video id}{title}{description}; thumbnail in assets/lesson2/video/yt-<id>.jpg
\newcommand{\ytclip}[3]{%
  \begin{minipage}[t]{5.15cm}
    \centering
    \href{https://www.youtube.com/watch?v=#1}{%
      \begin{tikzpicture}
        \node[inner sep=0pt] (img) {\includegraphics[height=3.35cm]{assets/lesson2/video/yt-#1.jpg}};
        \draw[hair,line width=1pt] (img.south west) rectangle (img.north east);
        \fill[black,opacity=.55] (img.center) circle (0.42cm);
        \fill[white] ($(img.center)+(-0.13,-0.2)$) -- ($(img.center)+(-0.13,0.2)$) -- ($(img.center)+(0.21,0)$) -- cycle;
      \end{tikzpicture}}\par\vspace{0.35em}
    {\color{ink}\small\bfseries #2}\par\vspace{0.1em}
    {\color{muted}\scriptsize #3}
  \end{minipage}}

\newcommand{\kicker}[1]{}
\newcommand{\subt}[1]{{\color{muted}\small #1}\par\vspace{.65em}}
\newcommand{\tagline}[1]{\takeaway{#1}}
\newcommand{\citefoot}[1]{}
\newcommand{\boxtitle}[1]{{\color{ink}\bfseries #1}}
\tcbset{l2base/.style={enhanced,arc=2pt,outer arc=2pt,boxrule=0.55pt,colframe=hair,colback=white,left=2mm,right=2mm,top=1.5mm,bottom=1.5mm}}
\newcommand{\coursebox}[2]{\begin{tcolorbox}[l2base]\boxtitle{#1}\par\vspace{1mm}{\small #2}\end{tcolorbox}}
\newcommand{\sboxblue}[2]{\begin{tcolorbox}[l2base,colback=tint,colframe=ink!24]\boxtitle{#1}\par\vspace{1mm}{\small #2}\end{tcolorbox}}
\newcommand{\lightimg}[2]{\begin{tcolorbox}[l2base,left=0pt,right=0pt,top=0pt,bottom=0pt]\includegraphics[width=\linewidth,height=#2,keepaspectratio]{#1}\end{tcolorbox}}
\newcommand{\inlinecode}[1]{{\ttfamily\small\color{ink}#1}}
\newcommand{\lc}[1]{{\color{muted}\footnotesize #1}}
\tikzset{
  lbox/.style={draw=hair,rounded corners=2pt,fill=white,text width=2.35cm,align=center,inner sep=8pt,font=\small},
  lblue/.style={draw=ink,rounded corners=2pt,fill=tint,text width=2.35cm,align=center,inner sep=8pt,font=\small},
  larrow/.style={-{Stealth[length=5pt]},draw=ink,line width=.8pt}
}

\renewcommand{\demoframe}[3]{%
\begin{frame}{#1}
  \begin{columns}[T,onlytextwidth]
    \column{0.55\textwidth}
      \group{What we do}
      #2
    \column{0.40\textwidth}
      \group{What to look at}
      #3
  \end{columns}
\end{frame}}
